\begin{statementbox}
  两条互相垂直的椭圆切线，交点到中心的距离恒为 $\sqrt{a^2+b^2}$。

\begin{center}
\begin{minipage}{\linewidth}
\centering
\textbf{垂直切线的交点在同一个圆上}\par\smallskip
\begin{tikzpicture}[scale=1.1, every node/.style={font=\small}, line cap=round]
\draw[->,gray] (-4,0)--(4.2,0) node[right] {$x$};
\draw[->,gray] (0,-3.8)--(0,4) node[above] {$y$};
\draw[CStatement,thick] (0,0) ellipse [x radius=3,y radius=2];
\node[below left] at (0,0) {$O$};
\draw[CExplanation,dashed,thick] (0,0) circle [radius={sqrt(13)}];

\draw[CProof,thick] (3,-.65)--(3,3.2) node[above] {$x=a$};
\draw[CProof,thick] (-.7,2)--(3.65,2) node[right] {$y=b$};
\draw[CExample,thick] (0,0)--(3,2);
\draw (2.76,2)--(2.76,1.76)--(3,1.76);
\fill (3,2) circle(2pt) node[above right] {$P(a,b)$};
\node[below right] at (3,0) {$a$}; \node[above left] at (0,2) {$b$};
\node[fill=white,inner sep=3pt] at (0,-4.2) {蒙日圆：$x^2+y^2=a^2+b^2$};

\end{tikzpicture}
\par 椭圆实线、轨迹圆虚线；图取 $a=3,b=2$，全图等比例。
\end{minipage}
\end{center}


  \textbf{\textcolor{CStatement}{条件与结论}}
  \begin{description}[style=nextline, leftmargin=2.8em, labelsep=0.5em, itemsep=0.45em]
    \item[\textbf{条件（椭圆）：}]
      已知椭圆 $C:\frac{x^{2}}{a^{2}}+\frac{y^{2}}{b^{2}}=1$，其中 $a>b>0$，中心为 $O(0,0)$。
    \item[\textbf{条件（两条切线）：}]
      从椭圆外一点 $P(x_0,y_0)$ 向椭圆作两条不同的切线 $l_1,l_2$。
    \item[\textbf{结论一（正向）：}]
      若 $l_1\perp l_2$，则
      \[
        \boxed{x_0^2+y_0^2=a^2+b^2}.
      \]
    \item[\textbf{结论二（反向）：}]
      若 $P(x_0,y_0)$ 满足 $x_0^2+y_0^2=a^2+b^2$，则 $P$ 一定在椭圆外，而且过 $P$ 的两条椭圆切线互相垂直。
    \item[\textbf{补充方法（切点弦与韦达定理）：}]
      设两切点为 $A(x_1,y_1)$、$B(x_2,y_2)$，则切点弦 $AB$ 的方程是
      \[
        \boxed{\frac{x_0x}{a^2}+\frac{y_0y}{b^2}=1}.
      \]
      联立椭圆求两组坐标乘积，再利用两条切线的一般式判断垂直：
      \[
        PA\perp PB\ \Longleftrightarrow\
        \boxed{\frac{x_1x_2}{a^4}+\frac{y_1y_2}{b^4}=0}.
      \]
      \textbf{提醒：}这里不必计算 $x_1+x_2$、$y_1+y_2$，也不必分别求出两条切线的斜率。
    \item[\textbf{轨迹（蒙日圆）：}]
      两条垂直切线的交点 $P$ 的轨迹是
      \[
        \boxed{x^2+y^2=a^2+b^2},\qquad OP=\sqrt{a^2+b^2}.
      \]
      圆心为椭圆中心 $O$，半径为 $\sqrt{a^2+b^2}$。
    \item[\textbf{延伸（当原曲线是圆）：}]
      若 $a=b=r>0$，原曲线为 $x^2+y^2=r^2$，对应的轨迹圆为 $x^2+y^2=2r^2$，半径是原半径的 $\sqrt2$ 倍。
  \end{description}

\end{statementbox}
